\documentclass[11pt]{article}
\usepackage[margin=1in]{geometry}
\usepackage{amsmath,amssymb,amsthm}
\theoremstyle{definition}
\newtheorem{theorem}{Theorem}[section]
\newtheorem{proposition}[theorem]{Proposition}
\title{Analysis 11: Log-Power Weights, Crossing Scale,\\
and Concentration for Weighted Rolling Validation}
\author{Aileen Tao}
\date{}
\begin{document}
\maketitle

\section{Setup}
\[
Y_i=\theta+\xi_i,\qquad \xi_i\stackrel{\mathrm{iid}}{\sim}N(0,1),
\qquad 0<|\theta|<1.
\]
Here $Y_i$ is the observation at time $i$, $\theta$ is the unknown mean,
and $\xi_i$ are independent standard Gaussian errors. For current
training-sample size $n\in\{1,2,\ldots\}$, the training sample mean is
\[
\bar Y_n=\frac1n\sum_{j=1}^nY_j.
\]
Model 0 predicts $\bar Y_n$; Model 1 predicts zero. Both use squared loss.

Write $R_k(n)$ for Model $k$'s expected squared prediction error on
an independent new observation, and $\Delta_R(n)=R_0(n)-R_1(n)$
for the population-risk difference. Then
\[
R_0(n)=1+\frac1n,\qquad R_1(n)=1+\theta^2,\qquad
\Delta_R(n)=\frac1n-\theta^2.
\]
Define $n_{\rm pop}$ as the first training size with $\Delta_R(n)<0$:
\[
n_{\rm pop}=\lfloor\theta^{-2}\rfloor+1\sim\theta^{-2}\qquad(|\theta|\to0).
\]
\textbf{Proof.} See Appendix~\ref{app:population}.

\section{Weighted rolling validation}
Rolling validation (RV) predicts $Y_i$ from the first $i-1$
observations. Let $w_i>0$ be the deterministic weight assigned to
validation time $i$, whose training size is $i-1$; the weight sequence
is independent of the data and $\theta$.
Define the expected weighted-RV loss difference (Model 0 minus Model 1) by
\[
\Delta_{{\rm RV},w}(n)
=\sum_{i=2}^{n+1}w_i\left(\frac1{i-1}-\theta^2\right).
\]
The quantity $r_n$ is the weighted average of the reciprocal training sizes:
\[
r_n=\frac{\sum_{i=2}^{n+1}w_i/(i-1)}{\sum_{i=2}^{n+1}w_i},
\qquad
\frac{\Delta_{{\rm RV},w}(n)}{\sum_{i=2}^{n+1}w_i}
=r_n-\theta^2.
\]
Define $n_{{\rm RV},w}$ as the first training size with
$r_n<\theta^2$, provided such an integer exists. Positivity of the
weights gives
\[
r_n\ge\frac1n,\qquad n_{{\rm RV},w}\ge n_{\rm pop}.
\]
The crossing inequality holds whenever the weighted crossing exists.
Also $r_n$ strictly decreases, so there is a unique first strict crossing
and the inequality persists afterward.
\textbf{Proof.} See Appendix~\ref{app:weighted}.

\section{Log-power weights and crossing scale}
Take
\[
w_i=[\log(i+1)]^p,\qquad p>0\ \text{fixed}.
\]
The exponent $p$ controls the increase in weight given to later
validation times. As $n\to\infty$, integral comparison gives
\[
\sum_{i=2}^{n+1}w_i\sim n(\log n)^p,\qquad
\sum_{i=2}^{n+1}\frac{w_i}{i-1}
\sim\frac{(\log n)^{p+1}}{p+1}.
\]
Therefore,
\[
r_n\sim\frac{\log n}{(p+1)n}.
\]
\textbf{Proof.} See Appendix~\ref{app:sums}.

\begin{theorem}[First strict log-power crossing]
For fixed $p>0$, the expected weighted-RV crossing exists for every
$0<|\theta|<1$, and as $|\theta|\to0$,
\[
n_{{\rm RV},w}\sim\frac{\theta^{-2}\log(\theta^{-2})}{p+1}
=\frac{2\theta^{-2}}{p+1}\log(|\theta|^{-1}),
\]
\[
\frac{n_{{\rm RV},w}}{n_{\rm pop}}
\sim\frac{\log(\theta^{-2})}{p+1}
=\frac{2}{p+1}\log(|\theta|^{-1}).
\]
\end{theorem}
\noindent Log-power weighting delays the expected RV crossing by a
logarithmic factor relative to the population crossing.
\textbf{Proof.} See Appendix~\ref{app:crossing}.

\section{Concentration on the log-weight crossing scale}
Throughout this section, $p>0$ is fixed and $|\theta|\to0$. Suppose
\begin{equation}\label{eq:scale}
\frac{n}{n_{{\rm RV},w}}\to c,\qquad 0<c<\infty,
\end{equation}
where $n$ is the current training size, $n_{{\rm RV},w}$ is the
expected weighted-RV crossing size, and $c$ is fixed.
Using
\[
n_{{\rm RV},w}\sim\frac{\theta^{-2}\log(\theta^{-2})}{p+1},
\]
we obtain
\[
n\sim\frac{c}{p+1}\theta^{-2}\log(\theta^{-2}).
\]
Multiplying by $\theta^2$ gives
\[
n\theta^2\sim\frac{c}{p+1}\log(\theta^{-2}).
\]
Taking logarithms instead gives
\[
\log n=\log(\theta^{-2})+\log\log(\theta^{-2})+O(1)
\sim\log(\theta^{-2}).
\]
The $\log\log(\theta^{-2})$ term is lower order.
Here $c<1$ is before the expected crossing, $c=1$ is on the crossing
scale, and $c>1$ is after the expected crossing.

\subsection{Realized weighted-RV score}
Define the realized weighted-RV loss difference
\[
\Xi_{n,w}=\sum_{i=2}^{n+1}w_i
\left\{(Y_i-\bar Y_{i-1})^2-Y_i^2\right\}.
\]
Its sign determines the selected model: a negative value favors
Model 0, a positive value favors Model 1, and zero is a tie.

\subsection{Mean on the log-weight scale}
The exact mean identity is
\[
\mathbb E\Xi_{n,w}
=\Delta_{{\rm RV},w}(n)
=\left(\sum_{i=2}^{n+1}w_i\right)(r_n-\theta^2).
\]
For log-power weights,
\[
\sum_{i=2}^{n+1}w_i\sim n(\log n)^p,\qquad
r_n\sim\frac{\log n}{(p+1)n}.
\]
Thus the two contributions satisfy
\[
\left(\sum_{i=2}^{n+1}w_i\right)r_n
\sim\frac{(\log n)^{p+1}}{p+1},
\]
\[
\theta^2\sum_{i=2}^{n+1}w_i\sim n\theta^2(\log n)^p.
\]
Using the scaling relations established above, for fixed $c\ne1$,
\[
\mathbb E\Xi_{n,w}\sim
\frac{1-c}{p+1}[\log(\theta^{-2})]^{p+1}.
\]
The mean is eventually positive for $c<1$ and negative for
$c>1$; at $c=1$ the leading terms cancel.
\textbf{Proof.} See Appendix~\ref{app:mean}.

\subsection{Variance}
Start from the centered realized score
\[
\Xi_{n,w}-\mathbb E\Xi_{n,w}.
\]
Expanding the squared-loss difference gives two uncorrelated random
components: a quadratic Gaussian part and the linear Gaussian term
\[
-2\theta\sum_{i=2}^{n+1}w_i\xi_i.
\]
Because these components are uncorrelated,
\[
\operatorname{Var}(\Xi_{n,w})
=\operatorname{Var}(\text{quadratic part})
+\operatorname{Var}\!\left(-2\theta\sum_{i=2}^{n+1}w_i\xi_i\right).
\]
The quadratic variance satisfies
\[
\operatorname{Var}(\text{quadratic part})\sim2(\log n)^{2p}.
\]
For the linear component, the square-weight asymptotic
\[
\sum_{i=2}^{n+1}w_i^2\sim n(\log n)^{2p}
\]
gives
\[
\operatorname{Var}\!\left(-2\theta\sum_{i=2}^{n+1}w_i\xi_i\right)
=4\theta^2\sum_{i=2}^{n+1}w_i^2
\sim4n\theta^2(\log n)^{2p}.
\]
Using the scaling relations established above puts both variances
on the same scale:
\[
\operatorname{Var}(\text{linear part})
\sim\frac{4c}{p+1}[\log(\theta^{-2})]^{2p+1},
\]
\[
\operatorname{Var}(\text{quadratic part})
\sim2[\log(\theta^{-2})]^{2p}.
\]
The linear variance is larger by one logarithmic factor, with power
$2p+1$ rather than $2p$, so it dominates. Therefore,
\[
\operatorname{Var}(\Xi_{n,w})
\sim\frac{4c}{p+1}[\log(\theta^{-2})]^{2p+1}.
\]
\textbf{Proof.} See Appendix~\ref{app:variance}.

\section{Main concentration theorem}
\begin{theorem}[Concentration away from the crossing boundary]
For fixed $p>0$ and $c\in(0,\infty)\setminus\{1\}$,
\[
\frac{\operatorname{Var}(\Xi_{n,w})}{\{\mathbb E\Xi_{n,w}\}^2}
\sim\frac{4c(p+1)}{(1-c)^2\log(\theta^{-2})}\to0.
\]
Thus weighted RV selects Model 1 with probability tending to one for
$c<1$, and Model 0 with probability tending to one for $c>1$.
At $c=1$, the leading mean vanishes.
\end{theorem}
\noindent\textbf{Proof.} See Appendix~\ref{app:ratio}.

\appendix
\setlength{\abovedisplayskip}{7pt plus 2pt minus 2pt}
\setlength{\belowdisplayskip}{7pt plus 2pt minus 2pt}
\section*{Appendix}
\section{Population risk and weighted-RV formulas}
\subsection{Population risk}\label{app:population}
Let $Y_{\rm new}=\theta+\xi_{\rm new}$ be independent of the training
sample, with $\xi_{\rm new}\sim N(0,1)$. Then
$Y_{\rm new}-\bar Y_n=\xi_{\rm new}-n^{-1}\sum_{j=1}^n\xi_j$
has mean zero and variance $1+1/n$, whereas
$\mathbb EY_{\rm new}^2=1+\theta^2$. Subtracting the risks gives
$\Delta_R(n)=1/n-\theta^2$, negative exactly when $n>\theta^{-2}$.
The first such integer is $\lfloor\theta^{-2}\rfloor+1
=\theta^{-2}+O(1)$; equality of risks is a tie.

\subsection{Weighted expectation, monotonicity, and crossing}\label{app:weighted}
Independence of $Y_i$ and its training sample gives
\[
\mathbb E\{(Y_i-\bar Y_{i-1})^2-Y_i^2\}
=\frac1{i-1}-\theta^2.
\]
Linearity of expectation proves the stated weighted mean, without
requiring independent validation losses. Dividing by positive total
weight gives $r_n-\theta^2$. Each reciprocal is at least $1/n$,
so $r_n\ge1/n$ and a strict weighted crossing cannot precede
$n_{\rm pop}$. Moreover,
\[
r_{n+1}=
\frac{(\sum_{i=2}^{n+1}w_i)r_n+w_{n+2}/(n+1)}
{\sum_{i=2}^{n+1}w_i+w_{n+2}}<r_n,
\]
because $r_n\ge1/n>1/(n+1)$. Since $r_1=1>\theta^2$, if
$r_n\to0$ there is a first strict crossing, and the inequality
$r_n<\theta^2$ persists thereafter.

\section{Log-power weight sums and crossing scale}
\subsection{Weight-sum asymptotics}\label{app:sums}
All limits in this subsection take $n\to\infty$ with fixed $p>0$.
For fixed $0<\delta<1$, weights at $\delta n\le i\le n+1$ equal
$(\log n)^p(1+o(1))$ uniformly. Bounding the full sum above by
$n[\log(n+2)]^p$ and below by this last fraction gives
\[
1-\delta\le\liminf_{n\to\infty}
\frac{\sum_{i=2}^{n+1}w_i}{n(\log n)^p}
\le\limsup_{n\to\infty}
\frac{\sum_{i=2}^{n+1}w_i}{n(\log n)^p}\le1.
\]
Let $\delta\downarrow0$. Replacing $p$ by $2p$ proves
\[
\sum_{i=2}^{n+1}w_i^2\sim n(\log n)^{2p}.
\]
With $m=i-1$, the reciprocal-weight
summand is asymptotic to $(\log m)^p/m$. This comparison sum
diverges and its summand is eventually decreasing. Thus
\[
\sum_{i=2}^{n+1}\frac{w_i}{i-1}
\sim\sum_{m=2}^n\frac{(\log m)^p}{m}
\sim\int_2^n\frac{(\log x)^p}{x}\,dx
\sim\frac{(\log n)^{p+1}}{p+1}.
\]
Dividing the sums gives the formula for $r_n$, in particular $r_n\to0$.

\subsection{First strict crossing}\label{app:crossing}
Fix $0<\varepsilon<1$ and set
\[
n_-=\left\lfloor
\frac{(1-\varepsilon)\theta^{-2}\log(\theta^{-2})}{p+1}
\right\rfloor,\qquad
n_+=\left\lceil
\frac{(1+\varepsilon)\theta^{-2}\log(\theta^{-2})}{p+1}
\right\rceil.
\]
Then
\[
\log n_\pm=\log(\theta^{-2})+\log\log(\theta^{-2})
-\log(p+1)+\log(1\pm\varepsilon)+o(1),
\]
so $\log n_\pm/\log(\theta^{-2})\to1$ and
\[
\frac{r_{n_\pm}}{\theta^2}
\sim\frac{\log n_\pm}{(p+1)n_\pm\theta^2}
\longrightarrow\frac1{1\pm\varepsilon}.
\]
Since $r_n\to0$, Appendix~\ref{app:weighted} ensures existence for
every admissible $\theta$. For sufficiently small $|\theta|$, the
limits above and monotonicity give $n_-<n_{{\rm RV},w}\le n_+$.
Rounding changes each endpoint by at most one. Divide by
$\theta^{-2}\log(\theta^{-2})/(p+1)$ and let
$\varepsilon\downarrow0$ to prove
$n_{{\rm RV},w}\sim\theta^{-2}\log(\theta^{-2})/(p+1)$.
The identity $\log(\theta^{-2})=2\log(|\theta|^{-1})$ and
$n_{\rm pop}\sim\theta^{-2}$ give the crossing formulas.

\section{Mean, variance, and concentration on the log-weight scale}
\subsection{Mean}\label{app:mean}
Appendix~\ref{app:sums} supplies deterministic remainders
$o(\log n)$ and $o(n)$, depending only on $n$ and fixed $p$, such that
\[
\mathbb E\Xi_{n,w}
=(\log n)^p\left\{\frac{\log n}{p+1}+o(\log n)
-\theta^2[n+o(n)]\right\}.
\]
This proves the general expansion without a restriction on $n\theta^2$.
Under~\eqref{eq:scale}, $n\sim
c\theta^{-2}\log(\theta^{-2})/(p+1)\to\infty$, and
\[
n\theta^2\sim\frac{c}{p+1}\log(\theta^{-2}),\qquad
\log n=\log(\theta^{-2})+\log\log(\theta^{-2})
+\log\frac{c}{p+1}+o(1).
\]
In particular $\log n\sim\log(\theta^{-2})$ and
$n\theta^2=O(\log n)$, so $\theta^2o(n)=o(\log n)$. Hence
\[
\mathbb E\Xi_{n,w}
=(\log n)^p\left\{\frac{\log n}{p+1}-n\theta^2\right\}
+o\bigl((\log n)^{p+1}\bigr).
\]
This additive expansion holds also at $c=1$. For
$c\ne1$, its leading coefficient is nonzero, giving the
mean equivalent in Section 4.2. At $c=1$ it proves only
$\mathbb E\Xi_{n,w}=o([\log(\theta^{-2})]^{p+1})$.

\subsection{Variance}\label{app:variance}
In this proof write $S_m=\sum_{j=1}^m\xi_j$ and
$f(x)=[\log(x+2)]^p$, so $w_{m+1}=f(m)$. Expanding one loss gives
\[
(Y_{m+1}-\bar Y_m)^2-Y_{m+1}^2
=\frac{S_m^2}{m^2}-\frac{2\xi_{m+1}S_m}{m}
-2\theta\xi_{m+1}-\theta^2.
\]
Let $Q_{n,w}$ denote the centered quadratic Gaussian part of the score.
The centered score is $Q_{n,w}$ plus
$-2\theta\sum_{i=2}^{n+1}w_i\xi_i$, where
\[
Q_{n,w}=\sum_{m=1}^nf(m)
\left\{\frac{S_m^2-m}{m^2}-\frac{2\xi_{m+1}S_m}{m}\right\}.
\]
Each centered quadratic monomial has zero covariance with each
$\xi_i$: Gaussian third moments vanish and independent centered
factors have zero expectation. Thus the quadratic and linear terms
are uncorrelated. Collect the centered-square and pair-product
coefficients as
\[
a_t=\sum_{m=t}^n\frac{f(m)}{m^2}\quad(1\le t\le n),\qquad
b_l=\sum_{m=l}^n\frac{f(m)}{m^2}-\frac{f(l-1)}{l-1}
\quad(2\le l\le n+1),
\]
with an empty sum equal to zero. Then
\[
Q_{n,w}=\sum_{t=1}^na_t(\xi_t^2-1)
+2\sum_{l=2}^{n+1}b_l\sum_{t=1}^{l-1}\xi_t\xi_l.
\]
Distinct monomials in this expression are uncorrelated; centered
squares have variance $2$ and pair products have variance $1$.
Consequently,
\begin{equation}\label{eq:exactvar}
\operatorname{Var}(\Xi_{n,w})
=2\sum_{t=1}^na_t^2+4\sum_{l=2}^{n+1}(l-1)b_l^2
+4\theta^2\sum_{i=2}^{n+1}w_i^2.
\end{equation}

We now estimate the quadratic part as $n\to\infty$ with fixed $p>0$.
Integral comparison gives
$a_t\le C f(t)/t$ for a constant $C$ depending only on $p$, hence
$\sum_ta_t^2=O(1)$ since $\sum_tf(t)^2/t^2<\infty$.
Define the tail and residual coefficients
\[
T_n=\sum_{m=n+1}^{\infty}\frac{f(m)}{m^2},\qquad
B_l=\sum_{m=l}^{\infty}\frac{f(m)}{m^2}-\frac{f(l-1)}{l-1}.
\]
Integral comparison and integration by parts yield
\[
T_n\sim\frac{f(n)}n,\qquad
|B_l|\le C\left\{\frac{(\log l)^{p-1}}l
+\frac{(\log l)^p}{l^2}\right\}\quad(l\ge3).
\]
Indeed, the sum--integral error is $O(f(l)/l^2)$, and
\[
\int_l^\infty\frac{f(x)}{x^2}\,dx
=\frac{f(l)}l+\int_l^\infty\frac{f'(x)}x\,dx,\qquad
\int_l^\infty\frac{f'(x)}x\,dx
=O\!\left(\frac{(\log l)^{p-1}}l\right).
\]
For completeness, $f'(x)/x=O((\log x)^{p-1}/x^2)$, and for any
fixed real $q$,
\[
\int_l^\infty\frac{(\log x)^q}{x^2}\,dx
=\frac{(\log l)^q}{l}
\int_0^\infty e^{-v}\left(1+\frac v{\log l}\right)^q\,dv
\sim\frac{(\log l)^q}{l}.
\]
Dominated convergence applies: for $q\le0$ the factor is at most
$1$, and for $q>0$ it is bounded by $(1+v)^q$ when $\log l\ge1$.
Taking $q=p-1$ proves the derivative bound. It also gives
$\int_l^\infty f(x)/x^2\,dx=f(l)/l\{1+O(1/\log l)\}$,
which proves the tail and $a_t$ estimates. The sum--integral bound
applies for large $l$ because $f(x)/x^2$ is eventually decreasing.
The shift from $f(l)/l$ to
$f(l-1)/(l-1)$ is $O(f(l)/l^2)$.
Thus, writing $D_n=\sum_{l=2}^n(l-1)B_l^2$,
\[
D_n=\begin{cases}
O(1),&0<p<1/2,\\
O(\log\log n),&p=1/2,\\
O((\log n)^{2p-1}),&p>1/2,
\end{cases}
\qquad D_n=o(f(n)^2).
\]
These bounds follow by integrating $(\log x)^{2p-2}/x$.
For $l\le n$, $b_l=B_l-T_n$, whereas $b_{n+1}=-f(n)/n$.
Cauchy--Schwarz now gives
\[
\begin{split}
\sum_{l=2}^n(l-1)b_l^2
&=T_n^2\frac{n(n-1)}2
+O\!\left(D_n+T_n n\sqrt{D_n}\right)\\
&\sim\frac{f(n)^2}2.
\end{split}
\]
The endpoint contributes $n b_{n+1}^2=f(n)^2/n=o(f(n)^2)$.
Since $f(n)\to\infty$, the bounded $a_t$ contribution is also
negligible. Therefore
\[
\operatorname{Var}(Q_{n,w})\sim2f(n)^2\sim2(\log n)^{2p}.
\]
Finally, Appendix~\ref{app:sums} gives
$\sum_iw_i^2\sim n(\log n)^{2p}$. On scale~\eqref{eq:scale},
$n\theta^2\to\infty$, so the linear term in~\eqref{eq:exactvar}
dominates the quadratic term. Substituting the limits from
Appendix~\ref{app:mean} proves the full variance equivalent in Section 4,
also for $c=1$.

\subsection{Ratio and selection}\label{app:ratio}
For $c\ne1$, the mean is eventually nonzero. Dividing the
variance equivalent by the squared mean equivalent gives
\[
\frac{\operatorname{Var}(\Xi_{n,w})}{\{\mathbb E\Xi_{n,w}\}^2}
\sim
\frac{\frac{4c}{p+1}[\log(\theta^{-2})]^{2p+1}}
{\frac{(1-c)^2}{(p+1)^2}[\log(\theta^{-2})]^{2p+2}}
=\frac{4c(p+1)}{(1-c)^2\log(\theta^{-2})}\to0.
\]
Writing $\mu=\mathbb E\Xi_{n,w}\ne0$ in this proof, a wrong sign or
tie implies $|\Xi_{n,w}-\mu|\ge|\mu|$. Therefore
\[
\begin{cases}
\Pr(\Xi_{n,w}\le0),&\mu>0,\\
\Pr(\Xi_{n,w}\ge0),&\mu<0
\end{cases}
\ \le\Pr\{|\Xi_{n,w}-\mu|\ge|\mu|\}
\le\frac{\operatorname{Var}(\Xi_{n,w})}{\mu^2}\to0.
\]
The sign of the mean from Appendix~\ref{app:mean} then shows that weighted
RV selects Model 1 with probability tending to one for $c<1$,
and Model 0 for $c>1$. At $c=1$, only
$\mathbb E\Xi_{n,w}=o([\log(\theta^{-2})]^{p+1})$ follows;
no ratio limit is asserted there.
\end{document}
